% généré par scripts/06_textes.py
\dosmeta{Année de classement : 368 (directe : nom consulaire lu). Références : ICUR I 479 = ILCV 3650 = AE 1892, 31. Rattachement selon l’EDB : \emph{Incertae originis} (le volume I des ICVR réunit les inscriptions dont la provenance est perdue). Conservation selon l’EDB : \emph{s. Maria in Cappella}. Support : Plaque de marbre. 60 × 29 cm (h. × l.). Inscription gravée. Fiche EDB : Filippo Piazzolla, première version 2014. Notice du catalogue p.~\pageref{not:126}.}
\txtbloc{Texte transmis (EDCS, état 2022)}{\lnum{1}M(a)erentes parentes fecerunt quod\par
\lnum{2}ipse debuit facere infelices parentes\par
\lnum{3}dolo debere et hunc tumulum fecerunt\par
\lnum{4}qui natus est Fl(avio) Sergio et Fl(avio) Nigriniano conss(ulibus)\par
\lnum{5}III Idus Aug(ustas) die Saturnis depositus in hac\par
\lnum{6}domo aeterna dd(ominis) nn(ostris) Valentiniano et Valente\par
\lnum{7}Aug(ustis) iterum conss(ulibus) VII Idus Dece⸢m⸣b(res) die Saturnis\par
\lnum{8}qui vixit ann(os) XVIIII m(enses) III d(ies) XXVI\par
\lnum{9}in pace}
\txtbloc{Texte de l’EDB (EDB28428)}{\lnum{1}meren̂tes parentes fecerunt quod\par
\lnum{2}ipse debuit facere infelices parentes\par
\lnum{3}dolo debere et hunc tumulum fecerun̂t\par
\lnum{4}qui natus est Fl(avio) Sergio et Fl(avio) Nigriniano cons(ulibu)s\par
\lnum{5}III idus aug(ustas) die Saturnis depositus in hac\par
\lnum{6}domo aeterna dd(ominis) nn(ostris) Valentiniano et Valente\par
\lnum{7}aug(ustis) iterum cons(ulibu)s VII idus decenb(res) die Saturnis\par
\lnum{8}qui vixit ann(is) XVIIII m(ensibus) III d(iebus) XXVI\par
\lnum{9}in pace}
